\section{Analiza bezpieczeństwa: TLS i ograniczone role}
\label{sec:bezpieczenstwo}

%<*sec-metoda>
Rozdział mierzy wpływ dwóch mechanizmów bezpieczeństwa na operacje CRUD
i~na rezerwację miejsc przy współbieżności. Pierwszy to szyfrowanie połączenia
między aplikacją a~bazą (TLS), drugi to praca aplikacji na koncie
z~minimalnymi uprawnieniami zamiast na koncie administratora. Pomiar
wykonaliśmy na osobnym stosie (nakładka \texttt{docker-compose.security.yml}),
na zbiorze 1 mln biletów, z~indeksami.

\begin{table}[H]
\centering
\small
\caption{Konfiguracja bezpieczeństwa w~eksperymencie.}
\label{tab:sec-konfig}
\begin{tabular}{>{\raggedright\arraybackslash}p{1.9cm}>{\raggedright\arraybackslash}p{6.9cm}>{\raggedright\arraybackslash}p{5.9cm}}
\toprule
SZBD & konto aplikacyjne (profil \texttt{role}) & TLS (profil \texttt{tls}) \\
\midrule
PostgreSQL & rola \texttt{tb\_app}: \texttt{SELECT}, \texttt{INSERT}, \texttt{UPDATE}, \texttt{DELETE} przez \texttt{ALTER DEFAULT PRIVILEGES}; brak DDL i~własności tabel & \texttt{ssl=on}, klient \texttt{sslmode=verify-full} z~własnym CA \\
MariaDB & użytkownik \texttt{tb\_app}: te same cztery uprawnienia na \texttt{ticketbench.*} & \texttt{--ssl-cert}, klient weryfikuje CA (\texttt{tls=tb}) \\
MongoDB & uwierzytelnianie (SCRAM, plik klucza replica setu), rola z~akcjami \texttt{find}, \texttt{insert}, \texttt{update}, \texttt{remove}; bez \texttt{readWrite}, które pozwala usuwać kolekcje i~indeksy & \texttt{tlsMode=allowTLS}, klient \texttt{tls=true} z~CA \\
Valkey & użytkownik ACL \texttt{tbapp}: tylko prefiksy kluczy modelu (\texttt{w:*}, \texttt{r:*}, \ldots), bez kategorii \texttt{@admin} i~\texttt{@dangerous} (\texttt{FLUSHDB}, \texttt{CONFIG}, \texttt{KEYS}) & osobny port TLS 6380 \\
\bottomrule
\end{tabular}
\end{table}

Dane ładuje i~indeksy tworzy administrator, tak jak przy wdrożeniu
produkcyjnym. Następnie te same scenariusze i~test współbieżności (1 i~100
klientów, 500 spornych miejsc) przechodzą przez cztery profile połączenia:
\texttt{base} (administrator, bez TLS), \texttt{role}, \texttt{tls}
i~\texttt{tls\_role}. Przed pomiarem wykonujemy nieliczony przebieg
rozgrzewkowy wszystkich scenariuszy. Profile mierzymy w~kolejności
zrównoważonej (\texttt{base}, \texttt{role}, \texttt{tls}, \texttt{tls\_role},
a~potem odwrotnie), więc każdy profil ma dwa przebiegi po 5 powtórzeń
w~symetrycznych pozycjach i~liniowy dryf (nagrzewanie pamięci podręcznej,
przyrost danych) się znosi. Porównujemy mediany czasu po stronie serwisu,
który mierzy wywołanie sterownika, więc szyfrowanie i~deszyfrowanie TLS są
wliczone. Narzut grupy operacji to średnia geometryczna ilorazów
$t_{\mathrm{profil}}/t_{\mathrm{base}}$ po scenariuszach tej grupy.

Przed pomiarem sprawdziliśmy, że ograniczenia działają. Konto aplikacyjne
dostaje odmowę przy \texttt{DROP TABLE} w~PostgreSQL i~MariaDB, przy usunięciu
kolekcji w~MongoDB oraz przy \texttt{FLUSHDB} i~dostępie do klucza spoza
prefiksów w~Valkey (\texttt{NOPERM}). Widoki \texttt{pg\_stat\_ssl}, licznik
\texttt{Ssl\_accepts} i~lista klientów Valkey potwierdzają, że profile
\texttt{tls} i~\texttt{tls\_role} faktycznie szyfrują połączenia.
%</sec-metoda>

\begin{table}[H]
\centering
\small
\caption{Zmiana mediany czasu operacji względem profilu \texttt{base}
[\%], 1 mln biletów, z~indeksami. Wartość dodatnia oznacza spowolnienie.}
\label{tab:sec}
\tbinput{generated/tables/security.tex}
\end{table}

\begin{table}[H]
\centering
\small
\caption{Przepustowość rezerwacji [rez./s] w~teście współbieżności przy
1 i~100 klientach, strategie domyślne.}
\label{tab:sec-cont}
\resizebox{\linewidth}{!}{\tbinput{generated/tables/security_contention.tex}}
\end{table}

%<*sec-wyniki>
\paragraph{Punkt odniesienia.} Profil \texttt{base} mierzymy dwa razy:
na początku i~na końcu sekwencji. Różnica między tymi dwoma przebiegami tej
samej konfiguracji (kolumna ,,base A/B'' w~tab.~\ref{tab:sec}) wynosi
\secnoise{postgres} w~PostgreSQL, \secnoise{mariadb} w~MariaDB,
\secnoise{mongo} w~MongoDB i~\secnoise{valkey} w~Valkey. Dla operacji
trwających ułamek milisekundy to naturalna zmienność środowiska. Narzut
profilu, który się w~niej mieści, jest w~tym teście niemierzalny.

\paragraph{Ograniczone role nie kosztują.} Zmiana łącznego czasu po
przejściu na konto aplikacyjne wynosi \secpct{postgres}{role}{all}
w~PostgreSQL, \secpct{mariadb}{role}{all} w~MariaDB,
\secpct{mongo}{role}{all} w~MongoDB i~\secpct{valkey}{role}{all} w~Valkey.
Znaki są różne i~wszystkie wartości mieszczą się w~dryfie ,,base A/B''.
Sprawdzenie uprawnień to porównanie z~listą uprawnień roli. W~bazach
relacyjnych odbywa się raz przy rozpoczęciu wykonania instrukcji,
w~MongoDB raz na operację z~pamięci podręcznej sesji, a~w~Valkey raz na
polecenie (dopasowanie prefiksu klucza). Wobec kosztu dostępu do danych jest
to niezauważalne. Zasada najmniejszych uprawnień nie ma więc w~tym systemie
ceny wydajnościowej.

\paragraph{TLS przy pojedynczych operacjach.} Szyfrowanie połączenia zmienia
łączny czas o~\secpct{postgres}{tls}{all} (PostgreSQL),
\secpct{mariadb}{tls}{all} (MariaDB), \secpct{mongo}{tls}{all} (MongoDB)
i~\secpct{valkey}{tls}{all} (Valkey), również w~granicach dryfu. Pule
utrzymują połączenia otwarte, więc uzgadnianie sesji TLS (najdroższa część)
odbywa się raz na połączenie, a~nie na operację. Szyfrowanie symetryczne
kilkuset bajtów zapytania i~odpowiedzi kosztuje mikrosekundy. Profil
\texttt{tls\_role}, łączący oba mechanizmy, daje podobny obraz.

\paragraph{TLS przy wysokiej współbieżności.} Wyraźny efekt widać w~teście
rezerwacji (tab.~\ref{tab:sec-cont}). Valkey przy 100 klientach obsługuje
bez TLS \seccval{valkey}{base}{100}\,rez./s, a~z~TLS
\seccval{valkey}{tls}{100}\,rez./s. Valkey w~domyślnej konfiguracji (jeden
wątek I/O) sam szyfruje i~deszyfruje ruch wszystkich połączeń w~wątku, który
wykonuje polecenia, więc przy dziesiątkach tysięcy krótkich poleceń na
sekundę koszt TLS konkuruje bezpośrednio z~obsługą danych. Przy jednym
kliencie różnica jest mała (\seccval{valkey}{base}{1} wobec
\seccval{valkey}{tls}{1}\,rez./s), bo serwer nie jest wtedy wąskim gardłem.
W~bazach dyskowych zmiany przepustowości przy 100 klientach nie układają się
w~stały wzór: każda transakcja czeka na zapis dziennika, a~nie na
szyfrowanie.

\paragraph{Wniosek.} Konto o~minimalnych uprawnieniach i~TLS można włączyć
w~PostgreSQL, MariaDB i~MongoDB bez mierzalnej straty wydajności operacji
CRUD. W~Valkey TLS obniża przepustowość przy dużym ruchu. Trzeba to
uwzględnić przy planowaniu (wątki I/O, więcej instancji albo szyfrowanie na
poziomie sieci, np. w~prywatnej sieci klastra).
%</sec-wyniki>

